\chapter{QA Command Center Dashboard}

\section{Dashboard Mimarisi (Tek HTML + Chart.js)}
Dashboard tek dosya olarak \texttt{dashboard.html} içinde geliştirilmiştir. Chart.js grafikleri, Lucide ikonları ve \texttt{results.json} fetch ile çalışır. Serve kökü: \texttt{target/test-history/}.

\section{Genel Bakış ve Metrik Kartları}
Ana sayfada toplam test, geçen, kalan ve başarı oranı kartları bulunur. Her kart drill-down popup açarak ilgili test listesini gösterir.

\section{Test Geçmişi ve Modül Trendi}
``Genel Koşumlar'' ve ``Modül Bazlı'' sekmeleri ile geçmiş koşumlar timeline'da listelenir. Modül trendi grouped bar chart ile modül bazlı pass/fail eğilimini gösterir.

\section{Hata Analizi ve AI Notları}
Fail olan testler katman filtresine göre (UI / API / DB) listelenir. Her kartta hata mesajı, ekran görüntüsü, AI analiz notu ve ``Hatayı Raporla'' butonu yer alır.

\section{Modül Kapsamı ve OWASP Haritası}
Modül kontrol listesi risk tier rozetleri ve dairesel ilerleme halkaları ile gösterilir. OWASP test kategorileri modüllere eşlenmiştir.

\section{Raporlanan Hatalar (Master-Detail)}
Jira tarzı master-detail düzen: sol listede raporlanan hatalar, sağda detay paneli. Ekler bölümünde ekran görüntüsü ve hata logu indirme desteklenir. Kayıtlar \texttt{localStorage} ile persist edilir.

\section{results.json Veri Şeması}
Her koşum \texttt{runs[]} dizisinde; testler \texttt{caseId}, \texttt{status}, \texttt{module}, \texttt{screenshotPath}, \texttt{errorLogPath}, \texttt{aiAnalysis} alanlarıyla saklanır.

% TODO: Dashboard ekran görüntüsü ekle (Genel Bakış, Hata Analizi).
